\cleardoublepage

\section{绪论}

% ============================================================
% 本章约3000-4500字
% 根据《浙江大学本科生毕业论文（设计）编写规则》2.3.1.1：
%   "应包括论文的研究目的、流程和方法等。"
%
% 本章不设子章节，作为整体段落自然过渡。
% 内容顺序：研究背景 → 研究目的与意义 → 研究内容与技术路线 → 论文组织结构
% ============================================================

\par 拖延已成为当代社会普遍存在的自我调节问题。Steel 对涉及人格特质、动机与任务变量的 691 项实证研究进行系统元分析，将拖延界定为自主、非理性地推迟既定行动的行为模式，并从期望、价值、时间贴现与个体敏感性四个维度统一解释了拖延的发生机制\cite{steel2007nature}。后续综述显示，约 50\% 的大学生为慢性拖延者，其中约 20\% 的患者生活功能受损严重，达到重度水平\cite{klingsieck2013procrastination}。2025 年针对印度泰兰加纳邦医学生的横截面调查则报告，高达 77.6\% 的学生存在中度至高度学业拖延，其主要诱因依次为懒散、工作量超载以及被手机或睡眠分散注意力。拖延不仅影响学业与工作效率，更与职业倦怠、学业压力和心理健康问题显著相关。这些发现共同表明，拖延已远远超出个别现象的范畴，成为当代学习与工作生态中的结构性挑战。

\par 在拖延这一普遍现象之下，还存在一个特别值得关注的高风险亚群体——注意缺陷多动障碍患者。高校生群体中的 ADHD 症状流行率因诊断标准而异，可达2\%至12\%，患者常伴随学业表现下滑与心理健康共病\cite{nugent2014adhd}。从病理机制来看，ADHD 并非单纯的注意力缺失，而是涉及行为抑制、执行功能与情绪调节等多系统的神经发育性障碍\cite{barkley1997behavioral,adler2017structure}，其中情绪失调已被越来越多的研究确认为独立的核心维度\cite{solergutierrez2023emotion}。在时间感知方面，ADHD 个体并非无法监控时间，而是缺乏策略性的时间管理能力\cite{seesjarvi2025time}；而在学业影响方面，注意缺陷症状对长期学业成就的负面效应尤为突出\cite{henning2022adhd}。上述机制共同导致 ADHD 群体在任务启动、时间管理和情绪调节方面面临深层困境，也为后续数字化辅助工具的设计提供了明确的心理学靶点。

\par 然而，针对拖延与 ADHD 的既有干预手段在可及性、有效性与个性化适配性方面仍存在显著缺口。传统认知行为疗法是当前疗效最为确切的拖延干预方案，但受限于面对面形式和长期治疗师参与，时间与地理可及性不足。数字化认知行为疗法程序部分缓解了可及性问题，但用户完成率仅为34\%至36\%，用户黏性构成严峻挑战\cite{mutter2023studicare}。通用自我管理方法虽广泛应用，但缺乏对 ADHD 用户认知特点的理解，其刚性结构甚至可能加剧任务厌恶情绪。商业 ADHD 应用的系统综述更是揭示了商业市场与学术研究之间严重的知识鸿沟：极少数应用包含开发依据，几乎没有应用提供疗效的循证证据\cite{pasarelu2020attention}。这些局限共同指向了一个核心现实，即现有工具解决了记录任务的问题，但未能有效回答如何开始的问题。

\par 大语言模型及其智能体范式的演变为上述困境提供了全新的技术路径。2017 年 Transformer 架构的提出，通过自注意力机制实现全局依赖建模，为后续 GPT 系列、通义千问等所有现代大语言模型提供了底层架构基础\cite{vaswani2017attention}。Brown 等以 1750 亿参数训练得到的 GPT-3 首次系统性地展示了上下文学习能力，使模型无需梯度更新即可凭自然语言指令完成新任务\cite{brown2020gpt3}。Ouyang 等提出的 InstructGPT 通过人类反馈强化学习将模型与人类意图对齐，证明经过对齐的中等规模模型在用户偏好上可超越未经对齐的超大模型\cite{ouyang2022instructgpt}。DeepSeek、通义千问等国产开源模型在关键基准上已达到与国际旗舰同级的能力水平，使云端大语言模型应用接口的定价降至应用层可负担的工程成本。在架构层面，ReAct 智能体范式通过思考、行动与观察的交替循环，使大语言模型能够主动分析需求、调用工具并完成任务，突破了传统被动问答的封闭循环\cite{yao2022react}。Spring AI Alibaba 等工程框架进一步将 ReAct 范式、工具调用和检查点持久化等能力封装为开箱即用的工业级组件，使智能体开发在 Java 生态中具备工程化落地的成熟度\cite{spring_ai_alibaba}。

\par 从市场维度审视，数字疗法正借助人工智能技术突破传统疗法的固定化瓶颈。多家市场研究机构的数据显示，全球 ADHD 数字疗法市场规模在 2025 年约为 9 亿美元，预计到 2033 年将达到 25 至 26 亿美元\cite{adhdmarket2025}。中国儿童 ADHD 患病率约为 6.4\%，成人 ADHD 患病率估计为 2.5\% 至 3\%，但诊断率不足 20\%，大量患者尚未获得正式诊断，存在巨大的增量空间\cite{chinaburden2025}。与此同时，全球生产力应用市场规模在 2024 年约为 100 亿美元，用户已养成付费习惯，为新型 ADHD 辅助工具提供了现成的用户教育基础\cite{productivitymarket2025}。对话式人工智能作为低门槛交互范式，把所有功能折叠到自然语言输入的最小动作中，无需用户理解菜单层级或字段语义，恰好契合 ADHD 群体对低认知负担交互的需求。基于上述背景，本项目提出并实现了面向 ADHD 与拖延人群的 AI 智能体驱动辅助系统 Nudge，旨在通过自然语言交互降低使用门槛，通过主动任务分解化解启动难题，并通过情境感知与持久化任务清单减轻未完成任务对认知资源的持续占用。

\par 本项目的总体目标是设计并实现上述 Nudge 系统，将项目管理、人工智能对话与游戏化激励三大模块有机融合。围绕这一总体目标，本项目设定三项具体目标。其一，探索大语言模型智能体在认知辅助场景下的工程化落地方法，形成涵盖 ReAct 循环、工具调用、会话记忆与多模式切换的可复用技术架构。其二，验证人工智能智能体、项目管理、游戏化激励三元融合范式在 ADHD 与拖延群体中的可行性，使系统兼具结构深度与情感温度。其三，沉淀基于 Spring AI Alibaba 与国产开源大语言模型的本土化合规部署经验，为后续学术研究与产业落地提供可借鉴的工程参考。

\par 本项目的研究意义可从理论、实践和社会三个层面加以理解。在理论层面，本项目为大语言模型驱动的对话式智能体在心理健康领域提供了工程原型验证。Ni 与 Jia 的范围综述系统梳理了人工智能驱动的数字心理健康干预，将其映射至筛查分诊、治疗支持、远程监测、临床教育与人群预防五大场景，并明确指出大语言模型驱动的对话式智能体与共情沟通增强是当前最具发展潜力的方向之一\cite{ni2026scoping}。本项目则将 ReAct 智能体的思考、行动、观察循环范式落地到 ADHD 与拖延这一高度异质化、长期化且治愈率较低的具体场景中，将抽象的方法学路径转化为可被复现的应用案例。同时，本项目有助于填补成人 ADHD 辅助技术研究的明显缺口。Husain 对现有 ADHD 辅助技术的系统综述指出，绝大多数研究聚焦于儿童 ADHD，而成人群体在工作场所与高等教育场景中所面临的执行功能需求显著不同，相关辅助技术的研究与评估均明显不足\cite{husain2020investigating}。本项目以大学生与年轻知识工作者为目标用户，从积极支持视角而非纯治疗视角出发，结合三元融合的工程范式，为该缺口下的后续学术研究提供了一个新的参考样本。进一步地，本项目将情绪调节、执行功能与时间感知等心理学机制编码为智能体的提示词策略与任务分解逻辑，使心理学理论与对话式人工智能工程在同一系统内形成闭环验证，为该交叉方向上的实证研究提供可检验的工程基础。在实践层面，本项目为 ADHD 与拖延人群提供了一种低启动门槛、可日常使用的辅助工具，并通过基于 Spring Boot 与 Spring AI Alibaba 的完整架构沉淀了一条可复用的本土化工程路径，为后续开发者构建同类应用提供了具体可参照的架构样板。在社会层面，本项目以消费级路径切入，能够覆盖被传统医疗体系所忽视的、未确诊但具有真实困扰的更广泛拖延人群，扩大心理健康干预的可及性\cite{catala2020digital}。同时，本项目选择国产开源大语言模型以规避数据出境合规风险，从证据基础与技术合规两个层面响应了数字心理健康工具的循证性与公平性要求\cite{mohammed2023technology}。

\par 围绕上述目标，本项目的研究内容涵盖四个主要方面。第一，ADHD 辅助场景需求分析，识别任务启动困难、情绪调节失败以及未完成任务对认知资源的持续占用三方面核心痛点，并据此确定系统功能范围。第二，基础业务模块实现，包括用户管理、项目管理、任务与标签管理、数据持久化与缓存等基础支撑。第三，AI 智能体工作流实现，基于 ReAct 模式构建智能体核心，包含路由决策、工具调用与会话记忆。路由节点承担意图识别和复杂度判断两项关键分类决策，产生简单闲聊、简单构建、复杂闲聊和复杂构建四种处理分支，并包含自校验机制以降低路由误判风险。工具调用机制遵循最小权限原则和最小完备原则，分为查询类工具和行动类工具两大类。会话级上下文管理采用线程安全的哈希表以会话标识为键管理上下文实例，并引入检查点持久化机制支持对话中断后的状态恢复。第四，系统测试与工程化实践，包括功能测试、性能测试和持续集成与持续部署流水线建设。

\par 本项目的研究路线划分为四个递进阶段。阶段一为需求分析，通过文献调研和竞品分析识别 ADHD 用户核心需求，明确系统功能与非功能需求。阶段二为架构设计与技术选型，采用领域驱动设计的分层架构思想，将系统自上而下划分为接口层、业务层、领域层、基础设施层和数据库层五个层次，每层仅向下依赖以避免跨层调用和循环依赖。在此基础上，依据限界上下文原则将系统拆分为公共模块、用户模块、项目管理模块、智能体模块、安全模块和语音服务模块六个独立模块，各模块拥有独立的包结构和持久化实现，通过明确接口进行交互。阶段二还需要确定数据库 Schema、设计基于第七版 UUID 的主键策略和 TIMESTAMPTZ 时间戳规范，并完成各模块的接口契约定义。阶段三为核心模块开发，依次实现公共基础设施、用户模块、项目管理模块、AI 智能体模块、安全模块和语音服务模块，其中 AI 智能体模块基于 ReAct 模式实现路由决策、工具调用和会话记忆，是开发工作的核心重点。阶段四为测试与部署，编写功能测试用例并开展性能基准测试，搭建基于 GitHub Actions 的持续集成与持续部署流水线，采用多阶段 Docker 构建策略完成容器化部署，保证开发、测试与生产环境的一致性。

\par 除上述核心研究内容外，本项目的论文组织结构如下：第 2 章介绍背景调研与需求分析，从拖延与 ADHD 的理论机制、智能体技术成熟度、市场与竞品分析三个维度论证本项目的立项依据，并明确系统的功能需求与非功能需求；第 3 章介绍系统实现，详述系统架构设计、数据库设计、公共基础设施及各模块的详细设计与实现，涵盖用户模块、项目管理模块、AI 智能体模块、安全模块和语音服务模块；第 4 章介绍软件开发文档，包括 RESTful 接口设计规范、系统使用说明、功能测试结果、性能测试结果及持续集成与持续部署流水线设计；第 5 章为结论，总结核心研究成果，客观分析研究局限，并提出未来工作方向。
